%%%%%%%%%%%%%%%%%%%%%%%%%%%%%%%%%%%%%%%%%%%%%%%%%%%%%%%%%%%%%%%%%%%%%%%%%%%
%
% Generic template for TFC/TFM/TFG/Tesis
%
% By:
%  + Javier Macías-Guarasa.
%    Departamento de Electrónica
%    Universidad de Alcalá
%  + Roberto Barra-Chicote.
%    Departamento de Ingeniería Electrónica
%    Universidad Politécnica de Madrid
%
% By: + Javier Macías-Guarasa. Departamento de Electrónica Universidad de Alcalá + Roberto Barra-Chicote. Departamento de Ingeniería Electrónica Universidad Politécnica de Madrid
% 
% Based on original sources by Roberto Barra, Manuel Ocaña, Jesús Nuevo, Pedro Revenga, Fernando Herránz and Noelia Hernández. Thanks a lot to all of them, and to the many anonymous contributors found (thanks to google) that provided help in setting all this up.
%
% See also the additionalContributors.txt file to check the name of additional contributors to this work.
%
% If you think you can add pieces of relevant/useful examples, improvements, please contact us at (macias@depeca.uah.es)
%
% You can freely use this template and please contribute with comments or suggestions!!!
%
%%%%%%%%%%%%%%%%%%%%%%%%%%%%%%%%%%%%%%%%%%%%%%%%%%%%%%%%%%%%%%%%%%%%%%%%%%%

\chapter{Conclusiones y líneas futuras}
\label{cha:conc}

En esta última sección del documento se va a realizar lo siguiente. En primer lugar, se van a realizar una serie de comentarios,
a modo de conclusión, sobre el trabajo desarrollado y acerca de los resultados del experimento. Seguidamente, se van a describir
una serie de líneas de continuación para continuar y mejorar el presente trabajo.

\section{Conclusiones}
\label{sec:conc}

Este trabajo ha permitido implementar cuatro metaheurísticas clásicas para el Problema del Viajante desde cero, y evaluar
su comportamiento replicando, en la medida de lo posible, el diseño experimental de un estudio de referencia \cite{halim2019}. 
Los resultados obtenidos muestran una tendencia relativa consistente con dicho estudio: la búsqueda tabú se posiciona como la 
técnica más fiable en términos de calidad de solución, seguida por la colonia de hormigas, mientras que el algoritmo genético 
es, con diferencia, el más sensible tanto al tamaño de la instancia como a la configuración de sus hiperparámetros.

Precisamente esta sensibilidad sea, quizás, la conclusión más relevante del trabajo. No es casual que TS haya resultado ser
la técnica más consistente: su ajuste depende, esencialmente, de dos parámetros (\textit{tabu tenure} y número de iteraciones),
frente a los muchos más grados de libertad de GA (tamaño de población, probabilidades de cruce y mutación, tamaño de
torneo, número de élites), sin contar que estos operadores pueden, además, sustituirse por otros de naturaleza muy distinta
(cruce OX frente a PMX, mutación por inversión frente a intercambio, etc.). Un espacio de configuración más amplio implica
más formas de que una elección de hiperparámetros resulte inadecuada, y esto se hace especialmente visible al introducir
un criterio de parada por estancamiento común a las cuatro técnicas: un mismo valor de \texttt{patience} puede resultar
perfectamente razonable para una metaheurística y catastrófico para otra, como se ha comprobado con GA en las instancias
de mayor tamaño. En otras palabras, buena parte de la diferencia observada entre algoritmos no refleja necesariamente una
superioridad intrínseca de unas técnicas sobre otras, sino la distinta dificultad de encontrar una configuración de
hiperparámetros única que funcione razonablemente bien para todas las instancias evaluadas.

En cuanto al objetivo original del trabajo, replicar un estudio comparativo bajo un presupuesto computacional mucho más
reducido, los resultados respaldan la decisión tomada: aunque las métricas absolutas de calidad son peores que las del
estudio de referencia (cuyas ejecuciones llegaron a durar hasta 19 horas), la clasificación relativa entre algoritmos se
mantiene estable en gran medida. Esto sugiere que un presupuesto computacional ajustado permite obtener conclusiones
cualitativamente similares sin necesidad de una inversión de tiempo desproporcionada, lo cual es coherente con el propio
propósito de las metaheurísticas: renunciar a la garantía del óptimo global a cambio de tiempos de ejecución razonables.

\section{Líneas futuras}
\label{sec:fut-lin}

A partir del trabajo realizado, se pueden identificar varias líneas de continuación que podrían ampliar o mejorar los 
resultados obtenidos:

\begin{itemize}
  \item \textbf{Hiperparámetros dependientes de la escala de la instancia}. A lo largo de este trabajo, se ha identificado
  en repetidas ocasiones que ciertos hiperparámetros (la temperatura inicial de SA, la constante $Q$ de ACO, o el propio
  \texttt{patience}) dependen de la escala del problema, y un valor fijo no es igualmente adecuado para instancias de
  tamaños muy distintos. Una línea de trabajo natural consistiría en derivar estos valores automáticamente a partir de
  características de la propia instancia, como la distancia media entre ciudades o el número de nodos, en lugar de fijarlos
  previamente.

  \item \textbf{Metaheurísticas híbridas}. Habiendo implementado y validado las cuatro técnicas clásicas de forma
  independiente, un siguiente paso natural sería combinar sus mecanismos, siguiendo el ejemplo de trabajos como
  \cite{jingyan2012} o \cite{xu2023}. Dada la robustez observada en TS frente al crecimiento de la instancia, y la buena
  capacidad exploratoria de ACO, una hibridación entre ambas podría resultar especialmente prometedora.

  \item \textbf{Selección automática de metaheurística}. Como se ha comprobado, ninguna técnica resulta óptima en todos
  los escenarios. Empleando el marco de trabajo ya desarrollado, se podría explorar un enfoque de meta-aprendizaje, similar
  al de \cite{kanda2016}, que seleccione automáticamente la metaheurística más adecuada en función de las características
  de cada instancia.

  \item \textbf{Ampliación del experimento}. Aumentar el número de repeticiones permitiría aplicar pruebas de significancia 
  estadística (ANOVA, test de Tukey) similares a las del estudio de referencia, en lugar de limitarse a estadística puramente
  descriptiva. Asimismo, también sería interesante incorporar las instancias de mayor tamaño del estudio original (hasta 442 
  nodos) una vez se aborde la siguiente línea futura.

  \item \textbf{Optimización del coste computacional de la búsqueda tabú}. Como ya se apuntaba en el capítulo de
  implementación, TS recalcula el coste completo de la ruta para cada vecino evaluado. Implementar un cálculo incremental,
  que solo tenga en cuenta la variación de coste de las dos conexiones modificadas por cada movimiento \textit{2-opt},
  reduciría notablemente su coste por iteración, actualmente el más elevado de las cuatro técnicas.

  \item \textbf{Extensión a otras variantes del problema}. El trabajo se ha limitado al TSP simétrico. Sería interesante
  evaluar el comportamiento de estas mismas metaheurísticas frente a variantes como el TSP asimétrico, el TSP con ventanas
  de tiempo, o el TSP con múltiples viajantes, descritas en el capítulo~\ref{cha:desarrollo}, para comprobar si las
  conclusiones obtenidas se generalizan a otros contextos.

  \item \textbf{Aplicación a un caso real}. Por último, más allá de las instancias estandarizadas de TSPLIB, resultaría
  valioso aplicar el marco de trabajo desarrollado a un caso de uso real, como la planificación de rutas logísticas u otras
  aplicaciones ya mencionadas en la sección~\ref{sec:aplicaciones}.
\end{itemize}

%%% Local Variables:
%%% TeX-master: "../book"
%%% End:


